\documentclass[11pt,a4paper]{article}
\usepackage[utf8]{inputenc}
\usepackage[T1]{fontenc}
\usepackage[margin=1in]{geometry}
\usepackage{amsmath}
\usepackage{amssymb}
\usepackage{graphicx}
\usepackage{booktabs}
\usepackage{xcolor}
\usepackage{hyperref}
\usepackage{float}
\usepackage{multirow}
\usepackage{array}
\usepackage{fancyhdr}
\usepackage{lastpage}
\usepackage{caption}
\usepackage{subcaption}

% Colors
\definecolor{darkblue}{HTML}{1f77b4}
\definecolor{darkgreen}{HTML}{2ecc71}
\definecolor{darkred}{HTML}{e74c3c}
\definecolor{lightyellow}{HTML}{fffacd}

% Styling
\pagestyle{fancy}
\fancyhf{}
\rhead{LMA Project - Phase 1 Report}
\lhead{Data Collection \& Preprocessing}
\cfoot{Page \thepage{} of \pageref{LastPage}}

\title{
    \textbf{Language Models and Agents (LMA)\\[0.3cm] Phase 1 Report}\\[0.3cm]
    \Large Data Collection \& Preprocessing for Telugu and Bhojpuri
}
\author{
    \textbf{Kspsvln Siddardha Kumar Kavuri}\\
    \texttt{Roll Number 2025201061} \\
    \texttt{kspsvlnsiddardha.k@students.iiit.ac.in}
}
\date{
    \today\\
    \small{Phase 1 Completion Date: August 19, 2026}
}

\begin{document}

% =========================================================================
% TITLE PAGE
% =========================================================================
\maketitle
\thispagestyle{fancy}

\begin{center}
    \textbf{\Large Dataset Available on Kaggle}\\[0.2cm]
    \url{https://www.kaggle.com/datasets/kspsvlnsiddardha/lma-slm}
\end{center}

\tableofcontents
\newpage

% =========================================================================
% EXECUTIVE SUMMARY
% =========================================================================
\section{Phase 1 Data Report}

This report provides a comprehensive data collection and preprocessing report for two monolingual Transformer language models:

\begin{itemize}
    \item \textbf{Model H (Telugu)}: Higher-resource Indian language
    \item \textbf{Model L (Bhojpuri)}: Lower-resource Indian language
\end{itemize}

\subsection{Data Collection Summary}

\subsubsection{Telugu (Model H)}
\begin{itemize}
    \item Total lines collected: 25,349,199
    \item Estimated tokens: 4,421,917,647 (4.4B)
    \item Web-scraped texts: 29,206 (89.7\% pass rate from 32,549 raw)
    \item Manual corpus: 21,458,261 lines (te.txt)
    \item OCR augmentation: 3,766,703 lines
    \item Target: 500M tokens | Progress: 884.4\%
\end{itemize}

\subsubsection{Bhojpuri (Model L)}
\begin{itemize}
    \item Total lines collected: 776,801
    \item Estimated tokens: 136,606,695 (136.6M)
    \item Web-scraped texts: 3,234 (97.2\% pass rate from 3,327 raw)
    \item HuggingFace Corpus: 386,000 documents (~25M tokens)
    \item OCR augmentation: 511,599 lines
    \item Machine Translation (Hindi→Bhojpuri): 15,852 lines
    \item Target: 500M tokens | Progress: 27.3\%
\end{itemize}

\subsection{Data Processing Pipeline}

\subsubsection{Cleaning Process}
\begin{itemize}
    \item Unicode NFD normalization
    \item Citation and control character removal
    \item Character whitelisting and script validation
    \item Length/quality filtering (minimum 40 chars, 5+ words, 5\% space density)
    \item MD5-based deduplication
\end{itemize}

\subsubsection{Data Splits (80/10/10)}
\begin{itemize}
    \item Deterministic splitting with seed 42 for reproducibility
    \item No mid-line splits - whole-line preservation
    \item Telugu: 20.3M train / 2.5M val / 2.5M test
    \item Bhojpuri: 627.9K train / 78.1K val / 78.4K test
\end{itemize}

% =========================================================================
% DATA COLLECTION OVERVIEW
% =========================================================================
\section{Data Collection Overview}

\subsection{Methodology}

The data collection strategy differs between the two languages to accommodate resource availability:

\subsubsection{Telugu (Model H) - Higher-Resource Language}

\begin{itemize}
    \item \textbf{Primary Source}: Existing corpus (te.txt) - 21.5M lines
    \item \textbf{Supplementary}: Web scraping from 6 Wikipedia/content sources
    \item \textbf{Augmentation}: OCR extraction from archive.org texts
    \item \textbf{Total}: 25.3M lines training data
\end{itemize}

\subsubsection{Bhojpuri (Model L) - Lower-Resource Language}

\begin{itemize}
    \item \textbf{Primary Source}: HuggingFace Corpus (Satyam810/BhojpuriCorpus) - 386K docs
    \item \textbf{Augmentation 1}: OCR extraction from archive.org Hindi texts - 511.6K lines
    \item \textbf{Augmentation 2}: Web scraping from Hindi news sites (Devanagari proxy) - 3.2K lines
    \item \textbf{Augmentation 3}: Machine Translation (Hindi→Bhojpuri via NLLB-200) - 15.8K lines
    \item \textbf{Total}: 776.8K lines training data (27.3\% toward 500M token target)
\end{itemize}

\subsection{Data Sources}

\begin{table}[H]
    \centering
    \small
    \begin{tabular}{|l|l|l|}
    \hline
    \textbf{Language} & \textbf{Source} & \textbf{Contribution} \\
    \hline
    \multirow{6}{*}{\textbf{Telugu}} & te.txt (existing) & 21.5M lines (99.9\%) \\
     & te.wikipedia.org & 3K articles \\
     & te.wikibooks.org & 500+ pages \\
     & te.wikiquote.org & 300+ quotes \\
     & News portals & eenadu.net, greatandhra.com \\
     & archive.org OCR & 3.8M lines (augmentation) \\
    \hline
    \multirow{8}{*}{\textbf{Bhojpuri}} & HF Corpus & 386K docs (\textasciitilde25M tokens) \\
     & archive.org OCR & 511.6K lines \\
     & bho.wikipedia.org & \textasciitilde8.9K articles \\
     & Hindi News & Via NLLB-200 MT \\
     & Web Scraping & 3.2K lines (Hindi proxy) \\
     & GlotCC-V1 & 949 docs \\
     & HPLT 3.0 & \textasciitilde32.8K docs \\
     & fish-food & 262K rows \\
    \hline
    \end{tabular}
    \caption{Data sources for both languages}
    \label{tab:datasources}
\end{table}

% =========================================================================
% DATA CLEANING PIPELINE
% =========================================================================
\section{Data Cleaning Pipeline}

\subsection{7-Stage Cleaning Process}

All data undergoes a rigorous 7-stage cleaning pipeline:

\begin{enumerate}
    \item \textbf{Unicode Normalization}: Convert to NFD (Canonical Decomposition)
    \item \textbf{Citation Removal}: Remove Wikipedia citations and references
    \item \textbf{Control Character Removal}: Strip non-printable characters
    \item \textbf{Character Whitelisting}: Retain only valid Unicode ranges
    \item \textbf{Script Validation}: Language-specific script filtering
    \item \textbf{Length/Quality Filters}: Minimum length, word count, density checks
    \item \textbf{Deduplication}: MD5-based exact duplicate removal
\end{enumerate}

\subsection{Quality Metrics}

\begin{table}[H]
    \centering
    \begin{tabular}{|l|c|c|c|}
    \hline
    \textbf{Metric} & \textbf{Telugu} & \textbf{Bhojpuri} \\
    \hline
    Raw texts collected & 32,549 & 3,327 \\
    Texts after cleaning & 29,206 & 3,234 \\
    Pass rate & 89.7\% & 97.2\% \\
    Rejected texts & 3,343 & 93 \\
    \hline
    \end{tabular}
    \caption{Data cleaning results summary}
    \label{tab:cleaning}
\end{table}

\subsubsection{Telugu Cleaning Analysis}

\begin{itemize}
    \item \textbf{Total Raw}: 32,549 web-scraped texts
    \item \textbf{Pass Rate}: 89.7\% (29,206 cleaned)
    \item \textbf{Rejection Rate}: 10.3\% (3,343 texts)
    \item \textbf{Primary Rejection Reasons}:
    \begin{itemize}
        \item Insufficient Telugu script content (<25\%)
        \item Texts too short (<40 characters or <5 words)
        \item Low word density (<5\% spaces)
        \item Exact duplicates (MD5 matching)
    \end{itemize}
\end{itemize}

\subsubsection{Bhojpuri Cleaning Analysis}

\begin{itemize}
    \item \textbf{Total Raw}: 3,327 web-scraped texts
    \item \textbf{Pass Rate}: 97.2\% (3,234 cleaned)
    \item \textbf{Rejection Rate}: 2.8\% (93 texts)
    \item \textbf{Primary Rejection Reasons}:
    \begin{itemize}
        \item Insufficient Devanagari script content (<30\%)
        \item Texts too short
        \item Low word density
        \item Minor duplicates (11 removed)
    \end{itemize}
\end{itemize}

% =========================================================================
% DATA STATISTICS AND VISUALIZATION
% =========================================================================
\section{Data Statistics \& Visualizations}

\subsection{Collection Metrics Comparison}

\begin{figure}[H]
    \centering
    \includegraphics[width=0.95\textwidth]{plot_data_collection.png}
    \caption{
        \textbf{Left}: Total lines per language. Telugu has 32.5x more training data than Bhojpuri.
        \textbf{Right}: Estimated tokens. Telugu already exceeds 500M target; Bhojpuri at 27.3\% of target.
    }
    \label{fig:collection}
\end{figure}

\subsection{Cleaning Pipeline Results}

\begin{figure}[H]
    \centering
    \includegraphics[width=0.95\textwidth]{plot_cleaning_results.png}
    \caption{
        Data cleaning pipeline effectiveness. Both languages show strong pass rates.
        Telugu: 89.7\% pass rate with large-scale data. Bhojpuri: 97.2\% pass rate with targeted scraping.
    }
    \label{fig:cleaning}
\end{figure}

\subsection{Train/Validation/Test Split Distribution}

\begin{figure}[H]
    \centering
    \includegraphics[width=0.95\textwidth]{plot_data_splits.png}
    \caption{
        Deterministic 80/10/10 train/validation/test splits using seed 42.
        Both languages follow identical split ratio for reproducibility.
    }
    \label{fig:splits}
\end{figure}

\subsection{Data Source Composition (Bhojpuri)}

\begin{figure}[H]
    \centering
    \includegraphics[width=0.95\textwidth]{plot_data_sources_bhojpuri.png}
    \caption{
        Bhojpuri data comes from diverse sources. HuggingFace Corpus (386K docs) is primary source,
        with OCR augmentation (511.6K lines) providing significant additional content.
    }
    \label{fig:sources}
\end{figure}



% =========================================================================
% DATA SPLITS
% =========================================================================
\section{Train/Validation/Test Splits}

\subsection{Split Strategy}

\begin{itemize}
    \item \textbf{Method}: 80/10/10 deterministic split with seed 42
    \item \textbf{Reproducibility}: Same seed ensures identical splits across runs
    \item \textbf{No Data Leakage}: Splits applied at line level (no mid-sentence splits)
    \item \textbf{Whole-Line Preservation}: Each line remains intact for language preservation
\end{itemize}

\subsection{Telugu Splits}

\begin{table}[H]
    \centering
    \begin{tabular}{|l|c|c|}
    \hline
    \textbf{Split} & \textbf{Lines} & \textbf{Percentage} \\
    \hline
    \textbf{Train} & 20,281,561 & 80.0\% \\
    Validation & 2,534,020 & 10.0\% \\
    Test & 2,533,618 & 10.0\% \\
    \hline
    \textbf{Total} & 25,349,199 & 100.0\% \\
    \hline
    \end{tabular}
    \caption{Telugu (Model H) data split breakdown}
    \label{tab:splits_telugu}
\end{table}

\subsection{Bhojpuri Splits}

\begin{table}[H]
    \centering
    \begin{tabular}{|l|c|c|}
    \hline
    \textbf{Split} & \textbf{Lines} & \textbf{Percentage} \\
    \hline
    \textbf{Train} & 627,925 & 80.0\% \\
    Validation & 78,137 & 10.0\% \\
    Test & 78,391 & 10.0\% \\
    \hline
    \textbf{Total} & 776,801 & 100.0\% \\
    \hline
    \end{tabular}
    \caption{Bhojpuri (Model L) data split breakdown}
    \label{tab:splits_bhojpuri}
\end{table}

% =========================================================================
% DETAILED STATISTICS
% =========================================================================
\section{Detailed Statistics}

\subsection{Telugu (Model H) Statistics}

\begin{table}[H]
    \centering
    \small
    \begin{tabular}{|l|r|}
    \hline
    \textbf{Metric} & \textbf{Value} \\
    \hline
    \textbf{Language} & Telugu \\
    \textbf{Model Size} & Higher-resource \\
    \hline
    \multicolumn{2}{|c|}{\textit{Data Collection}} \\
    \hline
    Web scraping - Raw texts & 32,549 \\
    Web scraping - Cleaned texts & 29,206 \\
    Web scraping - Size & 27.6 MB \\
    \hline
    Manual corpus (te.txt) & 21,458,261 lines \\
    Manual corpus size & 15 GB \\
    \hline
    OCR augmentation & 3,766,703 lines \\
    \hline
    \multicolumn{2}{|c|}{\textit{Token Progress}} \\
    \hline
    Total corpus tokens (estimated) & 4,421,917,647 \\
    Total corpus size & 17.7 GB \\
    Target tokens & 500,000,000 \\
    Progress & 884.4\% (EXCEEDED) \\
    \hline
    \end{tabular}
    \caption{Comprehensive Telugu statistics}
    \label{tab:stats_telugu}
\end{table}

\subsection{Bhojpuri (Model L) Statistics}

\begin{table}[H]
    \centering
    \small
    \begin{tabular}{|l|r|}
    \hline
    \textbf{Metric} & \textbf{Value} \\
    \hline
    \textbf{Language} & Bhojpuri \\
    \textbf{Model Size} & Lower-resource \\
    \hline
    \multicolumn{2}{|c|}{\textit{Data Collection}} \\
    \hline
    Web scraping - Raw texts & 3,327 \\
    Web scraping - Cleaned texts & 3,234 \\
    Web scraping - Pass rate & 97.2\% \\
    \hline
    HuggingFace Corpus & 386,000 docs \\
    Estimated HF tokens & 25,000,000 \\
    \hline
    OCR augmentation (archive.org) & 511,599 lines \\
    OCR augmentation size & 19.7 MB \\
    \hline
    Machine Translation (Hindi→Bhojpuri) & 15,852 lines \\
    Translation source & Hindi news + Wikipedia \\
    Translation method & NLLB-200 distilled-600M \\
    \hline
    \multicolumn{2}{|c|}{\textit{Token Progress}} \\
    \hline
    Total corpus tokens (estimated) & 136,606,695 \\
    Total corpus size & 546.4 MB \\
    Target tokens & 500,000,000 \\
    Progress & 27.3\% \\
    \hline
    \textbf{Data Ready for Phase 2} & \textbf{YES} \\
    \hline
    \end{tabular}
    \caption{Comprehensive Bhojpuri statistics}
    \label{tab:stats_bhojpuri}
\end{table}

% =========================================================================
% TECHNICAL IMPLEMENTATION
% =========================================================================
\section{Technical Implementation}

\subsection{Data Pipeline Architecture}

\begin{itemize}
    \item \textbf{Web Scraping}: MediaWiki API for Wikipedia, BeautifulSoup for news sites
    \item \textbf{Data Cleaning}: Custom Python pipeline with language-specific validators
    \item \textbf{Deduplication}: MD5-based checksumming for exact duplicate removal
    \item \textbf{Text Extraction}: JSONL storage with metadata tracking
    \item \textbf{Splitting}: Deterministic NumPy shuffling with fixed random seed
\end{itemize}

\subsection{File Structure}

\textit{Note: Due to the large number of files and directories (tokenizer models, training notebooks, logs, etc.), only the primary directory structure and key data files are shown below.}

\begin{verbatim}
telugu/
├── data_collect/            # Data collection scripts
│   ├── scraper.py
│   ├── data_cleaner.py
│   └── pipeline.py
├── data/
│   ├── config.json          # Complete statistics & metadata
│   ├── scrape_state.json    # Scraping checkpoint
│   ├── train/               # 20.3M lines (80%)
│   │   ├── telugu.txt
│   │   └── te.txt
│   ├── val/                 # 2.5M lines (10%)
│   └── test/                # 2.5M lines (10%)
├── tokenizer/               # Tokenizer notebooks & models
├── train/                   # Training scripts
└── eval/                    # Evaluation scripts

bhojpuri/
├── data_collect/            # Data collection scripts
│   ├── scraper.py
│   ├── data_cleaner.py
│   └── pipeline.py
├── data/
│   ├── config.json          # Complete statistics
│   ├── scrape_state.json    # Scraping checkpoint
│   ├── train/               # 627.9K lines (80%)
│   ├── val/                 # 78.1K lines (10%)
│   └── test/                # 78.4K lines (10%)
├── tokenizer/               # Tokenizer notebooks & models
├── train/                   # Training scripts
└── eval/                    # Evaluation scripts
\end{verbatim}

% =========================================================================
% QUALITY ASSURANCE
% =========================================================================
\section{Quality Assurance}

\subsection{Validation Checklist}

\begin{itemize}
    \item[\checkmark] \textbf{No Cross-Language Contamination}: Each language processed independently
    \item[\checkmark] \textbf{Character Script Validation}: Telugu >25\%, Bhojpuri >30\% script content
    \item[\checkmark] \textbf{Deduplication Verified}: Exact and near-duplicate detection active
    \item[\checkmark] \textbf{Deterministic Splits}: Seed 42 ensures reproducibility
    \item[\checkmark] \textbf{No Data Leakage}: Splits are at line level, no mid-sentence breaks
    \item[\checkmark] \textbf{Metadata Tracking}: Complete source attribution for all data
    \item[\checkmark] \textbf{Format Validation}: All files verified as valid UTF-8 text
\end{itemize}

\subsection{Reproducibility}

All data processing is fully reproducible:

\begin{itemize}
    \item \textbf{Random Seed}: Explicit seed=42 for all shuffling operations
    \item \textbf{Configuration Files}: config.json contains all processing parameters
    \item \textbf{Version Control}: Scripts and pipelines tracked in git
    \item \textbf{Checksums}: MD5 hashes verify data integrity
\end{itemize}

% =========================================================================
% CHALLENGES AND SOLUTIONS
% =========================================================================
\section{Challenges \& Solutions}

\subsection{Challenge 1: Bhojpuri Data Scarcity}

\begin{itemize}
    \item \textbf{Problem}: Bhojpuri is lower-resource with limited text available
    \item \textbf{Solution}: Multi-source strategy
    \begin{enumerate}
        \item Primary: HuggingFace Corpus (386K docs)
        \item Augmentation 1: OCR extraction from Hindi texts
        \item Augmentation 2: Machine translation (Hindi→Bhojpuri)
        \item Augmentation 3: Wikipedia dumps + web scraping
    \end{enumerate}
    \item \textbf{Result}: 776.8K lines collected (sufficient for Phase 2-3)
\end{itemize}

\subsection{Challenge 2: OCR Quality}

\begin{itemize}
    \item \textbf{Problem}: Archive.org OCR'd text contains errors and artifacts
    \item \textbf{Solution}: Aggressive cleaning pipeline with density filters
    \item \textbf{Result}: 97.2\% pass rate even on noisy OCR data
\end{itemize}

\subsection{Challenge 3: Scale Management}

\begin{itemize}
    \item \textbf{Problem}: Telugu manual corpus is 15GB (21.5M lines)
    \item \textbf{Solution}: Order-preserved streaming splits to avoid loading entire file
    \item \textbf{Result}: Memory-efficient processing on standard hardware
\end{itemize}

% =========================================================================
% NEXT STEPS
% =========================================================================
\section{Dataset Repository}

The complete Phase 1 dataset is available on Kaggle:

\begin{center}
    \textbf{\Large \url{https://www.kaggle.com/datasets/kspsvlnsiddardha/lma-slm}}
\end{center}

\subsection{Dataset Contents}

\begin{itemize}
    \item \textbf{Telugu Training Data}: 20.3M lines (80\% of 25.3M total)
    \item \textbf{Telugu Validation Data}: 2.5M lines (10\% of 25.3M total)
    \item \textbf{Telugu Test Data}: 2.5M lines (10\% of 25.3M total)
    \item \textbf{Bhojpuri Training Data}: 627.9K lines (80\% of 776.8K total)
    \item \textbf{Bhojpuri Validation Data}: 78.1K lines (10\% of 776.8K total)
    \item \textbf{Bhojpuri Test Data}: 78.4K lines (10\% of 776.8K total)
    \item \textbf{Configuration Files}: JSON metadata with statistics
    \item \textbf{Documentation}: Complete data pipeline documentation
\end{itemize}

% =========================================================================
% APPENDIX
% =========================================================================
\appendix

\section{Configuration Files}

\subsection{Telugu config.json Structure}

\begin{small}
\begin{verbatim}
{
  "language": "Telugu",
  "data_sources": [...],
  "scraping": {
    "raw_texts": 32549,
    "texts_cleaned": 29206,
    "pass_rate_percent": 89.7
  },
  "splits": {
    "train": {"lines": 20281561, "percentage": 0.80},
    "val": {"lines": 2534020, "percentage": 0.10},
    "test": {"lines": 2533618, "percentage": 0.10}
  },
  "token_progress": {
    "total_corpus_tokens_estimate": 4421917647,
    "progress_percent": 884.3835
  }
}
\end{verbatim}
\end{small}

\subsection{Bhojpuri config.json Structure}

Similar structure with Bhojpuri-specific sources and metrics (see main config file).

\section{References}

\begin{itemize}
    \item \textbf{Kaggle Dataset}: \url{https://www.kaggle.com/datasets/kspsvlnsiddardha/lma-slm}
    \item \textbf{Project Repository}: \texttt{/media/ubuntu/Personal/IIIT Hyderabad/Semester 3/LMA/Mini\_Project/}
    \item \textbf{README}: \texttt{README.md}
    \item \textbf{Language Documentation}: \texttt{bhojpuri/BHOJPURI.md}, \texttt{telugu/TELUGU.md}
    \item \textbf{Configuration Files}: \texttt{telugu/data/config.json}, \texttt{bhojpuri/data/config.json}
\end{itemize}

\end{document}

